\documentclass[a4paper]{article}
% Español (México) con XeLaTeX: fontspec para fuentes Unicode, polyglossia
% para la división silábica y los nombres del documento en la variante mexicana.
\usepackage{fontspec}
\usepackage{polyglossia}
\setdefaultlanguage[variant=mexican]{spanish}
% Los nombres chinos van como «pinyin (original)»: xeCJK da a los caracteres
% Han su propia fuente; sin ella XeLaTeX los omite con un aviso.
% FandolSong no cubre algunos caracteres de nombres (U+73FA); WenQuanYi Zen Hei sí.
\usepackage[AutoFallBack=true]{xeCJK}
\setCJKmainfont{FandolSong-Regular.otf}
\setCJKfallbackfamilyfont{\CJKrmdefault}{WenQuanYi Zen Hei}
% polyglossia no traduce los nombres de listings.
\AtBeginDocument{\providecommand{\lstlistingname}{}\renewcommand{\lstlistingname}{Listado}%
  \providecommand{\lstlistlistingname}{}\renewcommand{\lstlistlistingname}{Índice de listados}}
\usepackage{amsmath, amssymb}
\usepackage{graphicx}
\usepackage[margin=2.5cm]{geometry}
\usepackage[most]{tcolorbox}
\usepackage{etoolbox}
\usepackage{listings}
\usepackage{hyperref}
\usepackage{booktabs}
\usepackage{subcaption}
\usepackage{float}
\newtcolorbox{knowledgebox}[1]{enhanced, colback=blue!5!white, colframe=blue!75!black, colbacktitle=blue!75!black, coltitle=white, fonttitle=\bfseries, title={#1}, attach boxed title to top left={yshift=-2mm, xshift=2mm}, boxrule=1pt, sharp corners}
\newtcolorbox{importantbox}[1]{enhanced, colback=yellow!10!white, colframe=yellow!80!black, colbacktitle=yellow!80!black, coltitle=black, fonttitle=\bfseries, title={#1}, sharp corners}
\newtcolorbox{warningbox}[1]{enhanced, colback=red!5!white, colframe=red!75!black, colbacktitle=red!75!black, coltitle=white, fonttitle=\bfseries, title={#1}, sharp corners}
\lstset{language=Python, basicstyle=\ttfamily\small, keywordstyle=\color{blue}, stringstyle=\color{red!60!black}, commentstyle=\color{green!60!black}, breaklines=true, frame=single, numbers=left, numberstyle=\tiny\color{gray}, captionpos=b, extendedchars=true}

\newcommand{\notetitle}{Derivación del decaimiento a larga distancia de RoPE Attention}
\newcommand{\noteauthors}{Elaboradas a partir de materiales públicos del curso}
\newcommand{\notedate}{2025}
\newcommand{\videochannel}{Wudaokou Nashi (五道口纳什)}
\newcommand{\videopublishdate}{2025}
\newcommand{\videoduration}{aproximadamente 17 minutos}
\newcommand{\videourl}{}
\newcommand{\videocoverpath}{cover.jpg}
\newcommand{\repourl}{https://github.com/hqhq1025/ai-course-notes}


% Footer with repo info
\usepackage{fancyhdr}
\pagestyle{fancy}
\fancyhf{}
\fancyfoot[C]{\thepage}
\fancyfoot[R]{\footnotesize\color{gray}\href{\repourl}{AI Course Notes}}
\renewcommand{\headrulewidth}{0pt}
\renewcommand{\footrulewidth}{0pt}

\begin{document}
\begin{titlepage}
\centering
{\Large Notas del curso\par}\vspace{1.2cm}
{\huge\bfseries \notetitle\par}\vspace{0.8cm}
{\large \noteauthors\par}\vspace{0.3cm}
{\large \notedate\par}\vspace{1.2cm}
\ifdefempty{\videocoverpath}{}{\includegraphics[width=0.82\textwidth,height=0.45\textheight,keepaspectratio]{\videocoverpath}\par}
\vfill
\begin{tcolorbox}[width=0.9\textwidth, colback=black!2!white, colframe=black!60, sharp corners]
\textbf{Autor o canal del video}: \videochannel\par
\textbf{Fecha de publicación}: \videopublishdate\par
\textbf{Duración del video}: \videoduration\par
\ifdefempty{\videourl}{}{\textbf{Enlace del video}: \href{\videourl}{\nolinkurl{\videourl}}\par}\par
\textbf{Página del proyecto}: \href{\repourl}{\nolinkurl{\repourl}}\par
\end{tcolorbox}
\end{titlepage}
\tableofcontents
\newpage

\section{Introducción}

La entrega anterior presentó, desde una perspectiva geométrica, la frecuencia de rotación y las propiedades multiescala de RoPE. Esta entrega deriva a fondo la expresión matemática de cómo varía el Attention Score de RoPE en función de la distancia, para responder una pregunta central: \textbf{¿por qué RoPE puede codificar simultáneamente dependencias de corto y de largo alcance?}

\section{Decaimiento exponencial de $\theta_j$}

Recordemos el parámetro de frecuencia de RoPE:
$$
\theta_j = \text{base}^{-2j/d}
$$

Tomando el logaritmo de $\theta_j$:
$$
\log \theta_j = -\frac{2j}{d} \log(\text{base})
$$

Esta es una función \textbf{lineal} de $j$, con pendiente $-\frac{2\log(\text{base})}{d}$. Por lo tanto, $\theta_j$ en sí mismo decae de forma exponencial.

\section{Derivación del Attention Score}

\subsection{Expansión del Pre-Softmax Score}

Bajo RoPE, el attention score (pre-softmax) entre la Query de la posición $m$ y la Key de la posición $n$ es:

$$
s(m,n) = (R_m q)^T (R_n k) = q^T R_m^T R_n k = q^T R_{n-m} k
$$

Al expandir $R_{n-m}$ en cada subespacio 2D, la contribución de cada subespacio $j$ es proporcional a:

$$
\cos\big((m-n) \cdot \theta_j\big)
$$

\begin{importantbox}{Fórmula central}
El attention score es proporcional a la suma de las funciones coseno de todos los subespacios:
$$
s(m,n) \propto \sum_{j=0}^{d/2-1} \cos\big((m-n) \cdot \theta_j\big)
$$
Cada subespacio aporta una onda coseno de una frecuencia distinta, determinada por $\theta_j$.
\end{importantbox}

\subsection{El efecto de la distancia $\tau = m - n$}

Sea $\tau = m - n$ la distancia entre tokens; analicemos el comportamiento de $s(\tau)$:

\textbf{Distancia corta ($\tau$ pequeña)}:
\begin{itemize}
  \item En todos los subespacios, $\cos(\tau \cdot \theta_j) \approx 1$
  \item Los términos se suman con signo positivo y el score es muy alto
  \item El modelo atiende fuertemente a los tokens cercanos
\end{itemize}

\textbf{Distancia larga ($\tau$ grande)}:
\begin{itemize}
  \item Subespacios de alta frecuencia ($\theta_j$ grande, $j$ pequeña): $\tau \cdot \theta_j$ cambia bruscamente y el valor de $\cos$ oscila con rapidez
  \item Subespacios de baja frecuencia ($\theta_j$ pequeña, $j$ grande): $\tau \cdot \theta_j$ cambia lentamente y el valor de $\cos$ decrece con lentitud
  \item Las oscilaciones bruscas de la parte de alta frecuencia se \textbf{cancelan entre positivos y negativos} al sumarse
\end{itemize}

\subsection{La imagen física de la atenuación remota}

\begin{importantbox}{El mecanismo de atenuación remota de RoPE}
La atenuación a larga distancia de RoPE es el resultado conjunto de dos efectos:
\begin{enumerate}
  \item \textbf{Cancelación total de las dimensiones de alta frecuencia}: los subespacios del lado izquierdo rotan con rapidez y, a distancias grandes, se cancelan entre positivos y negativos (de forma similar al efecto de interferencia de ondas)
  \item \textbf{Descenso lento de las dimensiones de baja frecuencia}: los subespacios del lado derecho rotan con lentitud y conservan una asociación débil a distancia
\end{enumerate}
En conjunto: matemáticamente, RoPE \textbf{atiende principalmente a lo local}, pero \textbf{conserva la capacidad remota}.
\end{importantbox}

Esto se asemeja al \textbf{fenómeno de interferencia de ondas} de la física: ondas de distintas frecuencias se cancelan entre sí a distancia, pero ciertas frecuencias específicas (las bajas) siguen propagándose hasta lejos.

\subsection{Análisis de los casos extremos}

Para el subespacio más a la derecha ($j = d/2 - 1$):
$$
\theta_{d/2-1} = \text{base}^{-1} \approx \frac{1}{\text{base}}
$$

Si base = 10000, entonces $\theta \approx 10^{-4}$, una longitud de onda extremadamente larga. Aunque $\tau$ sea muy grande, el crecimiento de $\tau \cdot \theta$ sigue siendo muy lento y $\cos(\tau \cdot \theta) \approx 1$. Esta es la razón por la que se conserva la información remota.

\subsection{Resumen de la sección}

El attention score es la superposición de las ondas coseno de todos los subespacios. A larga distancia, los términos de alta frecuencia se cancelan y los de baja frecuencia se atenúan lentamente; juntos producen la propiedad de atenuación remota de RoPE. RoPE implementa de manera natural un patrón de atención de «atender a lo local + conservar lo remoto».

\section{Relación entre $\theta_j$ y la frecuencia de rotación}

\begin{table}[H]
\centering
\caption{Índice del subespacio y características de frecuencia}
\begin{tabular}{cccc}
\toprule
\textbf{Posición del subespacio} & \textbf{Valor de $j$} & \textbf{$\theta_j$} & \textbf{Característica} \\
\midrule
Izquierda (índice bajo) & Pequeño & Grande (alta frecuencia) & Rotación rápida, codifica distancias cortas \\
Derecha (índice alto) & Grande & Pequeño (baja frecuencia) & Rotación lenta, codifica distancias largas \\
\bottomrule
\end{tabular}
\end{table}

\begin{knowledgebox}{Codificación multiescala de RoPE}
RoPE tiene distintas frecuencias de rotación en distintas dimensiones. Si se considera el vector de embedding como un vector fila:
\begin{itemize}
  \item Cuanto más a la izquierda esté el sub-space: mayor es la frecuencia, más sensible a la posición relativa
  \item Cuanto más a la derecha esté el sub-space: menor es la frecuencia, menos sensible a la posición relativa
\end{itemize}
Esta característica multiescala es la clave del éxito de RoPE.
\end{knowledgebox}

\section{Síntesis y ampliación}

\begin{enumerate}
  \item Attention score es proporcional a $\sum_j \cos(\tau \cdot \theta_j)$, es decir, la superposición de ondas coseno de múltiples frecuencias
  \item A distancias largas, los términos de alta frecuencia se cancelan (efecto de interferencia), y los términos de baja frecuencia decaen lentamente
  \item RoPE realiza de manera natural el patrón de atención de «priorizar lo local y conservar lo remoto»
  \item En la próxima sesión se explorará el reconocimiento de patrones de los Attention Head y el fenómeno de Attention Sink
\end{enumerate}

\subsection{Lecturas adicionales}
\begin{itemize}
  \item Su et al., artículo original de «RoFormer»
  \item NIPS 2025 Best Paper: Gated Attention y Attention Sink
  \item Blog: «Understanding RoPE through the lens of signal processing»
\end{itemize}

\end{document}
